\subsection{APDU Message Mapping}
\label{sec:apdu-mapping}

The four logical messages of the PQ-Aliro Expedited Phase do not map directly
to single ISO-DEP APDUs.  Because post-quantum cryptographic objects far exceed
the 255-byte ISO/IEC~14443 APDU payload limit, every message must be
\emph{fragmented} into a chain of command or response APDUs before transmission.
On the Reader (command) side the implementation sets
\texttt{MAX\_CHUNK\_SIZE = 200\,B} for AUTH0 (\texttt{INS = 0x80}) and AUTH1
(\texttt{INS = 0x81}), and \texttt{MAX\_CHUNK\_SIZE = 255\,B} for LOADCERT
(\texttt{INS = 0xD1}).
Each intermediate chunk carries \texttt{P1 = 0x10} to signal that further data
follows; the final chunk sets \texttt{P1 = 0x00}.
On the User Device (response) side, \texttt{MAX\_APDU\_CHUNK = 240\,B}; the UD
appends status word \texttt{0x61\,\textit{xx}} to every non-final response chunk
and \texttt{0x9000} to the last.
The Reader retrieves subsequent chunks with a \texttt{GET RESPONSE}
(\texttt{INS = 0xC0}) command.

Fields whose encoded length exceeds 255 bytes are wrapped in BER-TLV
\emph{long-form} length encoding: a leading \texttt{0x82} byte followed by a
two-byte big-endian length, adding 4 bytes of framing overhead per such field.
This affects the ML-KEM-768 ephemeral public key (1\,184\,B $\to$ 1\,188\,B
on the wire), the KEM ciphertext (1\,088\,B $\to$ 1\,092\,B), and all
ML-DSA-65 signatures (3\,309\,B $\to$ 3\,313\,B).

Table~\ref{tab:apdu} enumerates the six APDU-layer messages of a single
PQ-Aliro transaction under the default ML-DSA-65 $\times$ ML-KEM-768
configuration.

\begin{table}[t]
\caption{PQ-Aliro APDU message mapping (ML-DSA-65 $\times$ ML-KEM-768, default configuration)}
\label{tab:apdu}
\centering
\small
\begin{tabular}{llclr}
\hline
\textbf{Message} & \textbf{Direction} & \textbf{INS} & \textbf{Payload Summary} & \textbf{Chunk} \\
\hline
AUTH0 Command    & Reader$\to$UD & \texttt{0x80} &
    KEM ephemeral $pk_e$ (1\,184\,B) $+$ \textit{trans\_id} (16\,B) $+$ \textit{reader\_id} (32\,B) &
    200\,B \\
AUTH0 Response   & UD$\to$Reader & ---           &
    KEM ciphertext (1\,088\,B) $+$ UD noise (32\,B) &
    240\,B \\
LOADCERT Command & Reader$\to$UD & \texttt{0xD1} &
    Profile0000 certificate ($\approx$5\,500\,B): $pk_{\mathrm{reader}}$ (1\,952\,B) $+$ signature (3\,309\,B) $+$ DN fields ($\approx$200\,B) &
    255\,B \\
LOADCERT Response & UD$\to$Reader & ---          &
    Status word only (\texttt{0x9000}) &
    --- \\
AUTH1 Command    & Reader$\to$UD & \texttt{0x81} &
    ML-DSA-65 transcript signature $\sigma_{\mathrm{reader}}$ (3\,309\,B) &
    200\,B \\
AUTH1 Response   & UD$\to$Reader & ---           &
    $\mathrm{AEAD}(pk_{\mathrm{ud}} \mathbin{\|} \sigma_{\mathrm{ud}} \mathbin{\|} \mathit{bitmap})$,
    plaintext 5\,273\,B $+$ GCM tag 16\,B $=$ 5\,289\,B &
    240\,B \\
\hline
\end{tabular}
\end{table}

Under this configuration the six messages collectively require approximately
75 APDUs and transfer $\approx$16.4\,KB of application-layer data, representing
a $\sim\!14\times$ increase over the $\approx$1.2\,KB of the classical ECC
instantiation.  The dominant contributors are the LOADCERT round
($\approx$5\,500\,B, 22 APDUs) and the AUTH1 Response
($\approx$5\,289\,B, 23 APDUs), both driven by the large size of ML-DSA-65
public keys and signatures.
